% 1. Cores e pacotes gráficos (devem vir PRIMEIRO)
\usepackage[usenames,dvipsnames,svgnames]{xcolor}
\usepackage{graphicx}
\usepackage{epsfig}
\usepackage{pgf,tikz}
\usetikzlibrary{calc,decorations.pathreplacing}
\usepackage{tkz-euclide}

% 2. Linguagem e codificação
\usepackage[utf8]{inputenc}
\usepackage[brazil]{babel}

% 3. PythonTeX
\usepackage{pythontex}

% 4. Matemática e fontes
\usepackage{amsfonts,mathtools,amsthm,amssymb}
\usepackage{stmaryrd}

% 5. Layout, tabelas e formatação
\usepackage{fancyhdr}
\usepackage{makeidx}
\usepackage{enumerate}
\usepackage{wrapfig}
\usepackage{xspace,setspace}
\usepackage[rm]{titlesec}
\usepackage[hmargin=1cm,vmargin=1.5cm]{geometry}
\usepackage{tabu}
\usepackage{multicol}
\usepackage{ccicons}
\usepackage{makecell}
\usepackage[nodayofweek]{datetime}

% 6. Caixas (tcolorbox)
\usepackage[many]{tcolorbox}
\tcbuselibrary{theorems}

% 7. Hyperref (DEVE vir por último entre os pacotes de formatação)
\usepackage{hyperref}
\hypersetup{
    colorlinks = true,
    urlcolor   = cyan,
    linkcolor  = red,
    citecolor  = blue
}

% 8. Licença
\usepackage[
type={CC},
modifier={by-nc},
version={4.0},
lang={brazilian},
]{doclicense}

% Configurações gerais
\renewcommand{\baselinestretch}{1.2}
\allowdisplaybreaks
\makeindex
%\nonstopmode

%----------------- Caixa de Teoremas (tcolorbox) --------------------------
\newtcbtheorem[number within=chapter]{theoremb}{Teorema}%
{
    enhanced,
    drop fuzzy shadow,
    colback=blue!5!white,
    colframe=blue!50!black,
    fonttitle=\bfseries,
    coltitle=white,
    description delimiters parenthesis,
    separator sign none
}{teo}

\newtcbtheorem[use counter from=theoremb]{exeb}{Exemplo}%
{
    enhanced,
    drop fuzzy shadow,
    colback=green!5!white,
    colframe=green!50!black,
    fonttitle=\bfseries,
    coltitle=white,
    description delimiters parenthesis,
    separator sign none
}{exe}

\newtcbtheorem[use counter from=theoremb]{exerb}{Exercício}%
{
    enhanced,
    drop fuzzy shadow,
    colback=red!5!white,
    colframe=red!50!black,
    fonttitle=\bfseries,
    coltitle=white,
    description delimiters parenthesis,
    separator sign none
}{exer}

%----------------- Teoremas Padrão --------------------------
\newtheorem{theorem}{Teorema}[chapter]
\newtheorem{lemma}[theorem]{Lema}
\newtheorem{corollary}[theorem]{Corolário}
\newtheorem{proposition}[theorem]{Proposição}

\theoremstyle{definition}
\newtheorem{definition}[theorem]{Definição}
\newtheorem{example}[theorem]{Exemplo}
\newtheorem{xca}[theorem]{Exercício}

\theoremstyle{remark}
\newtheorem{remark}[theorem]{Observação}

\numberwithin{section}{chapter}
\numberwithin{equation}{chapter}

%----------------- Novos Comandos --------------------------
\newcommand{\R}{\mathbb{R}}
\newcommand{\N}{\mathbb{N}}
\newcommand{\C}{\mathbb{C}}
\newcommand{\dt}[1]{{\color{blue}\textbf{#1}}}
\newcommand{\ep}{\varepsilon}
\newcommand{\nl}[1]{\|#1\|_{\mathcal{L}(X)}}
\newcommand{\dm}{d}
\newcommand{\vt}[1]{\overrightarrow{#1}}

\newcommand{\marcador}[2]{%
    \phantomsection%
    \addcontentsline{toc}{#1}{#2}%
    \noindent {\large \textbf{#2}}\par\vspace{1em}%
}

\newcommand{\n}[1]{\|#1\|}

\DeclareMathOperator{\rank}{\operatorname{rank}}
\DeclareMathOperator{\im}{Im}
\DeclareMathOperator{\id}{I}
\DeclarePairedDelimiterX{\inp}[2]{\langle}{\rangle}{#1, #2}


%----------------- Caixa de Sombra (Substituto do \shadowbox) --------------------------
\newtcbox{\myshadowbox}{
  enhanced,
  sharp corners,
  colback=white,
  colframe=black,
  boxrule=0.5pt,
  shadow={2pt}{-2pt}{0pt}{black}, % Reproduz a sombra clássica do \shadowbox
  boxsep=0pt,
  left=4pt, right=4pt, top=4pt, bottom=4pt,
  nobeforeafter,
  tcbox raise base
}

%----------------- Estilo do Capítulo --------------------------
\font\largefont= pzcmi scaled 3500
\font\numberfont= pzcmi scaled 3000

\titleformat{\chapter}[display]
  {\normalfont\Large}
  {\filright
    \rule[32pt]{.7\linewidth}{4pt}%
    \hspace{2pt}%
\raisebox{1cm}{	\myshadowbox{%
	  \begin{minipage}{.15\linewidth}
        \begin{center}
          \textsl{\large \chaptertitlename}\\[1ex]
          {\numberfont \thechapter}\\[1ex]
        \end{center}
      \end{minipage}%
    }}
  }
  {-10pt}
  {\filcenter\slshape\Huge}
  [\vspace{-2cm}\singlespacing\hfill\rule{.8\textwidth}{0.5pt}\\
   \vskip-2.8ex\hfill\rule{.7\textwidth}{4pt}\onehalfspacing]

\titlespacing{\chapter}{0pt}{*3}{*1}




\newtcbox{\destaque}{
enhanced,
    drop fuzzy shadow,
    colback=blue!5!white,
    colframe=blue!50!black
}

